% Solutions to the problems of Chapter 7 (Applications Beyond Autonomous Vehicles)
\begin{addedblock}
\begin{sol}{prob:ch07:1}
The only conflicting products are $m_1(\{N\})\,m_2(\{T\})=0.28$ and $m_1(\{N\})\,m_2(\{L,T\})=0.28$, so $K=0.56$. The surviving products are $m_1(\{N\})m_2(\Omega)=0.14\to\{N\}$, $m_1(\Omega)m_2(\{T\})=0.12\to\{T\}$, $m_1(\Omega)m_2(\{L,T\})=0.12\to\{L,T\}$ and $m_1(\Omega)m_2(\Omega)=0.06\to\Omega$; dividing by $1-K=0.44$ gives $m(\{N\})=0.318$, $m(\{T\})=0.273$, $m(\{L,T\})=0.273$, $m(\Omega)=0.136$. Hence $\mathrm{Bel}(T)=0.273$, $\mathrm{Pl}(T)=0.273+0.273+0.136=0.682$, and $\mathrm{BetP}(N)=0.318+0.136/3=0.364$, $\mathrm{BetP}(L)=0.273/2+0.136/3=0.182$, $\mathrm{BetP}(T)=0.273+0.273/2+0.136/3=0.455$. With $\alpha_1=0.5$ the AIS BBA becomes $m_1^{\alpha}(\{N\})=0.35$, $m_1^{\alpha}(\Omega)=0.65$; then $K=0.35\times0.8=0.28$, and the combined masses are $m(\{N\})=0.07/0.72=0.097$, $m(\{T\})=0.26/0.72=0.361$, $m(\{L,T\})=0.361$, $m(\Omega)=0.13/0.72=0.181$, so $\mathrm{BetP}(T)=0.361+0.181+0.060=0.602$. Discounting the less reliable source lowers the conflict and moves the decision towards the radar cue, which is the intended effect of reliability-aware discounting.
\end{sol}
\begin{sol}{prob:ch07:2}
Under $\mathcal{H}_0$ each $d_i^{2}$ is $\chi^{2}_{4}$ and consecutive innovations are independent, so the sum of $W=12$ of them is $\chi^{2}_{48}$. The one-sided $99\,\%$ threshold is $\chi^{2}_{48}(0.99)=73.7$ and the two-sided interval is $[\chi^{2}_{48}(0.005),\chi^{2}_{48}(0.995)]=[27.0,\,79.5]$; the Wilson--Hilferty formula with $z_{0.99}=2.326$ gives $48[1-2/432+2.326\sqrt{2/432}]^{3}=73.6$. Overlapping windows share $W-1$ terms, so successive statistics are strongly correlated and alarms arrive in runs; the alarm \emph{rate} is not $\alpha$ per window. The expected number of alarm samples is nevertheless bounded by $\alpha$ times the number of samples, i.e.\ at most $0.01\times360=3.6$ alarm samples per hour at a 10-s period (they will typically form fewer, longer runs).
\end{sol}
\begin{sol}{prob:ch07:3}
Write $\boldsymbol{\nu}=\mathbf{b}+\mathbf{e}$ with $\mathbf{e}\sim\mathcal{N}(\mathbf{0},\mathbf{S})$ and $\mathbf{z}=\mathbf{S}^{-1/2}\boldsymbol{\nu}\sim\mathcal{N}(\mathbf{S}^{-1/2}\mathbf{b},\mathbf{I})$; then $d^{2}=\mathbf{z}^{\mathsf T}\mathbf{z}$ is non-central $\chi^{2}_{n_y}$ with $\lambda=\|\mathbf{S}^{-1/2}\mathbf{b}\|^{2}=\mathbf{b}^{\mathsf T}\mathbf{S}^{-1}\mathbf{b}$, and the sum of $W$ independent such variables is non-central $\chi^{2}_{Wn_y}$ with non-centrality $W\lambda$ (non-centralities add). With $\sqrt{S_{11}}=31.3$\,m and a 50-m bias along one axis, $\lambda=(50/31.3)^{2}=2.55$. The per-step statistic has mean $n_y+\lambda=4.55$ against $\chi^{2}_{2}(0.99)=9.21$, so a typical sample does not fire; the windowed statistic has mean $Wn_y+W\lambda=24+30.6=54.6$ against $\chi^{2}_{24}(0.99)=43.0$, so it fires on average. The mean of the windowed statistic grows with $W$ while its threshold grows only like $Wn_y+O(\sqrt{W})$, hence the gain in sensitivity to small persistent biases, at the price of a delay of up to $W$ samples.
\end{sol}
\begin{sol}{prob:ch07:4}
$S(a_1)=0.376(0.90)+0.121(0.80)+0.215(0.40)+0.215(0.60)+0.074(0.70)=0.701$; similarly $S(a_2)=0.495$, $S(a_3)=0.584$, $S(a_4)=0.299$. Rescaling the other weights by $(1-w_1)/(1-0.376)$, the difference $S(a_1)-S(a_3)$ is affine in $w_1$: it equals $0.30\,w_1+(1-w_1)\,\delta$, where $\delta=[0.121(0.30)+0.215(0.10)+0.215(-0.35)+0.074(0.30)]/0.624=0.0083>0$ is the rescaled contribution of the remaining criteria. Both terms are non-negative for $w_1\in[0,1)$, so $a_3$ never overtakes $a_1$ by changing the proximity weight alone; the ranking can only be reversed through the innovation weight, as found in Example~\ref{ex:ch07-maritime-mcdm}.
\end{sol}
\begin{sol}{prob:ch07:5}
With $\varphi_n(x)=\sqrt{2/L}\sin(n\pi x/L)$, $H_{jn}=\frac{1}{L_g}\sqrt{2/L}\,[\sin(n\pi(x_j+5)/L)-\sin(n\pi(x_j-5)/L)]=\frac{2}{L_g}\sqrt{2/L}\cos(n\pi x_j/L)\sin(5n\pi/L)$. At $x_j=250$\,m the cosine factors are $\cos(\pi/4)$, $\cos(\pi/2)=0$, $\cos(3\pi/4)$ and $\cos(\pi)$, so numerically $H_{j1}=3.16\times10^{-3}\times\cos(\pi/4)\times\ldots$, and in particular $H_{j2}=0$: the channel sits at the strain node of mode~2 (an antinode of displacement), where $\varphi_2'$ vanishes. In the small-gauge approximation $H_{jn}\approx\varphi_n'(x_j)=\sqrt{2/L}\,(n\pi/L)\cos(n\pi x_j/L)$, a channel at $x_j=500$\,m has $\cos(n\pi/2)=0$ for every odd $n$, so it cannot observe modes 1, 3, 5, \ldots. For localisation the channel set must therefore avoid the common strain nodes of the retained modes (e.g.\ the mid-span and the quarter points) or, equivalently, use an irregular spacing so that the columns of $\mathbf{H}$ remain well conditioned.
\end{sol}
\begin{sol}{prob:ch07:6}
See \repopath{companion\_\allowbreak repo/\allowbreak exercises/\allowbreak ch07/\allowbreak solutions/\allowbreak prob06\_\allowbreak spoof\_\allowbreak sweep.py}. With the seed of the book example, the per-step test fires $570$, $420$, $350$ and $160$\,s after the onset for $a_s=0.0005$, $0.001$, $0.002$ and $0.004$\,m/s$^{2}$, at AIS offsets between $50$ and $120$\,m (two to four radar standard deviations). The window test is earlier for slow ramps ($500$\,s against $570$\,s at $a_s=0.0005$\,m/s$^{2}$ with $W=12$; $320$ against $350$\,s at $0.002$\,m/s$^{2}$) and later for the fast ramp ($210$ against $160$\,s at $0.004$\,m/s$^{2}$), the longest window $W=24$ being the slowest to react in every case: a long window integrates a small bias but also dilutes a sudden one. The Ljung--Box alarm fraction in the coloured regime is $0.00$, $0.04$, $0.09$ and $0.20$ for $\phi=0.5$, $0.7$, $0.85$ and $0.95$; it grows monotonically with $\phi$ and is negligible for $\phi\le0.5$.
\end{sol}
\begin{sol}{prob:ch07:7}
See \repopath{companion\_\allowbreak repo/\allowbreak exercises/\allowbreak ch07/\allowbreak solutions/\allowbreak prob07\_\allowbreak cable\_\allowbreak sweep.py}. The detection delay is one to three samples (longer with only three modes, whose strain response to a point load is weaker); the centroid error drops markedly from three to five modes (the modal resolution $L/(2n_m)$ shrinks) and then saturates at a few tens of metres, fluctuating with the channel geometry, with the largest errors for strikes near the ends, where the antisymmetric residual pattern is truncated by the boundary. A wrong coupling factor produces a residual excess confined to the affected channels for the whole record, whereas a mechanical event is transient and spreads over neighbouring channels.
\end{sol}
\end{addedblock}
